% Transcription — fonds Alexandre Grothendieck, shelfmark 145, batch 2 (pages 21-37).
% Established from the facsimile archives/batches/145/batch-02.pdf (a mirror of
% https://grothendieck.umontpellier.fr/145.pdf), per the skill
% transcribe-grothendieck. The folder is thirty-seven pages; this is the second
% and last batch (seventeen pages).
%
% All seventeen pages carry his hand and all seventeen are transcribed; none is
% skipped and none is summarised as another's text. Page 31 holds only four
% lines. Page 32 opens with his own title « (2) Le groupe S E^^ » (the 2
% circled); the opening heading of the batch is the edition's and says so.
% His own pagination: 9, 10 on pages 21, 22; 11, 12, 13, 14 on pages 24, 26,
% 28, 30; 15, 16, 17 on pages 32, 34, 36 — recorded as \note{p. N de
% l'auteur}; no number was seen on the pages between.
%
% What the batch is about. It continues batch 1's « Critique des notations »:
% the discrete extension E = E_{0,3} of Gamma = Gamma_{0,3} = S_3 x Z/2 by
% pi = pi_{0,3} = pi_1(P^1_C - {0,1,inf}), with l_inf l_1 l_0 = 1. Page 21
% (his p. 9): a system (10) for the quasi-square roots eps_i and rho (very
% uncertain), the isotropy of P in Gamma, the section (11) phi: S_3 -> Gamma,
% sigma -> (sigma, sg sigma), and its canonical lift (12) to E. Pages 22-25:
% the subgroup E_tau, « groupe cartographique de Malgoire », with (17) l_0 =
% tau_inf tau_1 etc.; the profinite completion E^ and the group E^^ of loop
% automorphisms of pi^, extension (18) of Gamma^^ = Autext_lac(pi^) by pi^;
% the commutant bold-Gamma of S_3 (Gamma n bold-Gamma = {1, tau}); u in E^^
% described by a multiplier p in Z^* and g_i in pi^ with u(l_i) =
% int(g_i^-1) l_i^p (19)-(22), composition (23), the extension SE^^! of E^^!
% by Z^{0,1,inf}, the 1-cocycle gamma_i (24), and the kernel pi^ x
% Z^{0,1,inf}; covariance of gamma_i under rho. Pages 26-28: partial
% splittings of E over Gamma — fixed points P, Pbar of rho, 2 = 0' of
% sigma_0, -1 = 1', 1/2 = inf'; the three arcs of U^tau; U^{tau sigma_1} =
% U - {1}; nothing new. Pages 28-30: a presentation of E, first attempted and
% struck (page 28 foot, page 29 top), then by sigma_i, rho, tau_i with the
% relations of page 29, then by rho, sigma_0, tau_0 alone: a quotient of
% (Z/2 x Z/2) * Z/3 by the single relation [rho, sigma_0][rho^-1, tau_0] = 1,
% and a symmetry sigma_0 <-> tau_0, rho <-> rho^-1 not compatible with the
% extension. Page 31: the extension SGamma of Gamma by Z^{0,1,inf} and SE =
% SGamma x_Gamma E still to be made explicit. Pages 32-37 (his pp. 15-17 and
% after): « Le groupe SE^^ » as the set of systems (p, sigma, g_0, g_1,
% g_inf) under condition (2), with (3)-(10) — the map to Aut_lac(pi^), the
% homomorphism i from pi^ x Z^^{0,1,inf}, the product law (8) and its
% non-explicit character, (9), (9'), (10); the partial section (11)-(12)
% phi: Gamma*_P -> SE, the subgroup g_i = 1 characterised by (13), and the
% « Exercice ! » that (14) l_inf^p l_1^p l_0^p = 1 forces p = 1; the Sylow
% 2-subgroups do not lift; int(rho~), int(sigma~_0) (15)-(16); the
% centralisers of rho~ (17)-(18) and of sigma~_0; the page ends on
% commutation in Gamma^^, i.e. modulo inner automorphisms.
%
% Reading hazards fixed once here, to agree with batch 1. His curly capital
% for the extension is set \mathcal{E} (read as a Sigma at first glance; batch
% 1 sets it \mathcal{E}); its completions \hat{\mathcal{E}},
% \hat{\hat{\mathcal{E}}}. The discrete pi_1 is a lower-case pi on these
% pages and is set \pi. His double-struck Gamma on page 23 is set
% \boldsymbol{\Gamma}. The exclamation mark in E^^!, Gamma^^!, Autext^!
% is his and is kept. « lac » under Autext / Aut is set as a subscript. The
% superscript -1 on g_i under int (pages 23-24) is added by him in a darker
% pencil; it is absent from pages 32-37 and is not supplied there.
%
% Readings to check first: the whole of system (10) on page 21 and the margin
% there; « \ill{} \ill{} » in the name of Malgoire's group on page 22; the
% oblique margin of page 23 (Gamma^^! and its « j'ignore si … »); the
% generator written over a hatched erasure on page 28, read rho' with
% « rho'^6 = 1, rho'^3 = tau » (the page seems to write rho^16, rho^13); the
% struck opening of page 29; « Question : … » struck at the foot of page 29;
% page 33, the insertion point of « mais ils s'explicitent ainsi »; the
% indices in (8) and (9); page 35, « \ill{} lieu de » and the bracket « quitte
% : corriger par un sigma~_i »; page 36, « 2-groupes de Sylow »; the struck
% lines and the oblique margin on page 37.
%
% Pass: Opus 5.5 (claude-opus-5-5), 2026-09-24 — first pass, unchecked against the pages by a human.
\documentclass[11pt,a4paper]{article}
% Le préambule commun à toutes les transcriptions du fonds Grothendieck.
%
% Il tient en une idée : **l'incertitude se compose, elle ne se résout pas.**
% Une transcription qui lisse un mot douteux a détruit la seule chose pour
% laquelle on la faisait — un lecteur doit pouvoir savoir, à chaque endroit,
% ce qui était sur le feuillet et ce qui a été ajouté. Les six macros qui
% suivent sont donc l'appareil critique, et non de la décoration.
%
% Les mêmes commandes sont reconnues par `scripts/render.mjs`, qui produit la
% vue de lecture du site. Une macro ajoutée ici doit l'être aussi là-bas,
% faute de quoi le rendu échouera bruyamment — ce qui est voulu : un rendu
% silencieusement faux est pire qu'un rendu absent.

\ProvidesPackage{grothendieck}[2026/08/08 Transcriptions du fonds Grothendieck]

\usepackage{amsmath,amssymb,amsthm}
\usepackage{mathrsfs}
\usepackage{stmaryrd}
% Les diagrammes commutatifs sont partout dans ce fonds : c'est la langue
% même de la géométrie algébrique de Grothendieck, et une transcription qui
% les rendrait en prose ne transcrirait rien.
\usepackage{tikz-cd}
% Français par défaut — c'est la langue des feuillets — mais l'anglais est
% chargé aussi : la traduction et le résumé sont des documents anglais, et la
% typographie française (espaces avant les deux-points) y serait une faute.
% Ces documents basculent par \selectlanguage{english} en tête de corps.
\usepackage[english,french]{babel}
\usepackage{fontspec}
\usepackage{geometry}
\geometry{a4paper,margin=25mm,marginparwidth=28mm}
\usepackage{marginnote}
\usepackage{xcolor}
% ulem, not soul. `soul`'s \st and \ul are built on a character-by-character
% scan that XeLaTeX's font machinery breaks: the document fails with an
% undefined control sequence rather than degrading. ulem does the same two jobs
% and is Unicode-engine safe. `normalem` stops it hijacking \emph, which the
% transcriptions use for Grothendieck's own emphasis.
\usepackage[normalem]{ulem}
\usepackage{hyperref}
\hypersetup{colorlinks=true,linkcolor=black,urlcolor=black,pdfborder={0 0 0}}

% --- Métadonnées du lot -----------------------------------------------------
% Reprises telles quelles par le script de rendu, qui les affiche en tête de
% la vue de lecture. Elles fixent ce que le fichier prétend couvrir : sans
% elles, une transcription flotte sans point d'attache dans un fonds de seize
% mille feuillets.

\newcommand{\folder}[1]{\gdef\grothFolder{#1}}
\newcommand{\batch}[1]{\gdef\grothBatch{#1}}
\newcommand{\leaves}[2]{\gdef\grothFirst{#1}\gdef\grothLast{#2}}
\newcommand{\foldertitle}[1]{\gdef\grothFoldertitle{#1}}
\gdef\grothFolder{?}\gdef\grothBatch{?}\gdef\grothFirst{?}\gdef\grothLast{?}\gdef\grothFoldertitle{}

% --- Le résumé --------------------------------------------------------------
%
% Il ouvre la lecture modernisée, sous le titre « Résumé », et il est composé
% en petit. Ce n'est pas un ornement : c'est le seul passage du document qui
% ne soit pas des mathématiques, et un lecteur doit pouvoir le reconnaître
% avant de l'avoir lu — puis le sauter s'il connaît déjà le sujet, ou s'y
% arrêter s'il ne le connaît pas. Le filet qui le suit marque la reprise du
% corps.

\newenvironment{resume}
  {\small\setlength{\parskip}{0.45em}}
  {\par\medskip\hrule\medskip}

% --- Mots-clés ---------------------------------------------------------------
%
% En anglais, en clôture du résumé de la lecture modernisée : le vocabulaire
% moderne sous lequel chercher ce que les feuillets construisent. Le manifeste
% du site (`npm run manifest`) les extrait pour étiqueter la cote entière —
% cette ligne est la source unique des tags, il n'y a pas de fichier à côté.
\newcommand{\keywords}[1]{\par\smallskip\noindent
  {\footnotesize\textsc{Keywords} --- \textit{#1}}\par}

% --- L'appareil critique ----------------------------------------------------

% Le numéro de feuillet, en marge. C'est l'ancre qui permet de régler un
% désaccord sur une formule en regardant *un* feuillet plutôt qu'une chemise
% de six cents. Le site s'en sert aussi pour tourner les pages du fac-similé
% quand on fait défiler la transcription.
\newcommand{\leaf}[1]{%
  \marginnote{\footnotesize\color{gray!70}\textsf{#1}}%
  \label{leaf:#1}\ignorespaces}

% Illisible : on ne devine pas. Un mot inventé qui se lit comme les autres est
% le pire résultat possible.
\newcommand{\ill}{\textcolor{red!60!black}{[\,\ldots\,]}}

% Lecture douteuse : le mot est donné, et signalé comme incertain.
\newcommand{\uncertain}[1]{\uline{#1}}

% Ajout de l'éditeur — une lettre manquante, un mot sous-entendu.
\newcommand{\add}[1]{\textcolor{blue!50!black}{[#1]}}

% Biffé par Grothendieck lui-même. Ce qu'il a rayé fait partie du document :
% une rature dit souvent où la pensée a changé de direction.
\newcommand{\struck}[1]{\sout{#1}}

% Note du transcripteur, sur l'état matériel du feuillet.
\newcommand{\note}[1]{\par\smallskip{\small\color{gray!60!black}\textit{[#1]}}\par\smallskip}

% Note marginale de Grothendieck, distincte du corps du texte.
\newcommand{\marginal}[1]{\par\smallskip\noindent\hspace{1em}%
  {\small\color{gray!60!black}#1}\par\smallskip}

% --- Titre ------------------------------------------------------------------

\AtBeginDocument{%
  \begin{center}
    {\large\bfseries Fonds Alexandre Grothendieck --- cote n\textsuperscript{o}~\grothFolder}\\[2pt]
    {\itshape\grothFoldertitle}\\[4pt]
    {\small feuillets \grothFirst--\grothLast\ (lot \grothBatch)}\\[2pt]
    {\footnotesize\color{gray!60!black}Transcription établie d'après le fac-similé numérisé
     par l'Université de Montpellier.\\
     Les lectures incertaines sont soulignées, les illisibles notées [\,\ldots\,],
     les ajouts entre crochets.}
  \end{center}
  \vspace{1em}}

\endinput


\folder{145}
\batch{2}
\pages{21}{37}
\dating{1981}
\watermark{Édition de démonstration}
\foldertitle{Printemps 81 (vieille rédaction) : notes manuscrites (s.d.)}

\begin{document}

\section{Critique des notations --- suite : l'extension $\mathcal{E}$ de $\Gamma_{0,3}$ par $\pi$}

\note{ce titre est de l'édition : le lot reprend au milieu du développement
commencé à la page 14 (lot 1), sans titre de sa main avant celui de la
page 32}

\page{21}
\note{p. 9 de l'auteur}

\marginal{il faut aussi écrire (au moins !) les relations
$\rho\,\varepsilon_0\,\rho^{-1} = \varepsilon_1$ etc., donc \uncertain{on a}
\ill{} \ill{} $(\varepsilon_0\rho)$\ill{}}

\[
(10)\quad
\left\{
\begin{aligned}
&\rho^{3} = 1\\
&\varepsilon_\infty^{2}\ \varepsilon_1^{2}\ \varepsilon_0^{2} = 1\\
&\varepsilon_1 = \varepsilon_0\,\varepsilon_\infty^{\uncertain{4}}\,\rho\,\varepsilon_{\uncertain{0}}^{2}\ \varepsilon_\infty^{2}\\
&\varepsilon_\infty = \varepsilon_1\,\varepsilon_0^{\uncertain{4}}\,\rho\,\varepsilon_{\uncertain{\infty}}^{2}\ \varepsilon_0^{2}\\
&\varepsilon_0 = \varepsilon_\infty\,\varepsilon_1^{\uncertain{4}}\,\rho\,\varepsilon_\infty^{2}\ \varepsilon_1^{2}
\end{aligned}
\right.
\]
\note{système lu tel quel et très incertain. Devant la première des trois
dernières lignes, un signe biffé ; entre le second $\varepsilon^{2}$ et le
dernier facteur de chacune de ces lignes, un petit trait vertical ondulé
surchargé, lu comme une rature et omis ; au-dessus de $\varepsilon_0$ de la
même ligne, une surcharge illisible. Les exposants lus « 4 » peuvent être
autre chose. La deuxième ligne pourrait vouloir dire
$\varepsilon_\infty^2 = \varepsilon_1^2 = \varepsilon_0^2 = 1$}

\uncertain{Chacun} \ill{} des $\varepsilon_i$, $\rho$ s'expriment en fonction
des $\varepsilon_i$, et on trouve comme \ill{} \uncertain{tantôt} \ill{}.

Quant à relever $\tau$ en un élément de $\mathcal{E}$, \uncertain{pour}
$\Gamma : \overline{P}$, on voit par des \uncertain{choix} privilégiés --- on
voit trois choix également naturels. Au lieu de les expliciter directement,
\uncertain{il est plus simple de} considérer le sous-groupe d'isotropie de
$P$ dans $\Gamma = \Gamma_{0,3} = \mathfrak{S}_3 \times \mathbb{Z}/2$, qui se
forme des couples $\dot\sigma, \dot\tau^{\varepsilon}$ tels que
$\operatorname{sg}(\dot\sigma) = \varepsilon$ ; \struck{\uncertain{donc}}
i.e.\ les éléments
\[
\dot\sigma^{*} = (\dot\sigma, \dot\tau^{\operatorname{sg}(\dot\sigma)}) \in \Gamma
\]
qui se relèvent canoniquement en des éléments ($\tilde\sigma$, \uncertain{fixant}
un $P$) de $\mathcal{E}$. Ces termes \uncertain{définissent} une section
\[
(11)\quad
\varphi : \mathfrak{S}_3 \longrightarrow \Gamma = \mathfrak{S}_3 \times \mathbb{Z}/2,
\qquad \dot\sigma \longmapsto (\dot\sigma, \operatorname{sg}(\dot\sigma))
\]
et un relèvement canonique de $\varphi$ en
\[
(12)\quad \sigma \longmapsto \tilde\sigma : \mathfrak{S}_3 \longrightarrow \mathcal{E} .
\]
On pose $\tilde\sigma_i = \tilde{\dot\sigma}_i$ (on a d'autre part
\struck{\ill{}} $\tilde\rho = \rho$)
\note{la parenthèse finale est tracée à gauche d'un passage biffé ; ce qui
suit « $\tilde\rho$ » est lu « $= \rho$ » sans certitude}

\page{22}
\note{p. 10 de l'auteur}

Le sous-groupe $\mathcal{E}_\tau$ de $\mathcal{E}$ image inverse de $\mathbb{Z}/2$
engendré par $\tau$, est le groupe cartographique « \ill{} \ill{} » de
Malgoire, ayant comme générateurs $\tau_0, \tau_1, \tau_\infty$
\struck{relation} avec
\[
(17)\quad
\left\{
\begin{aligned}
\struck{\tau_1\tau_\infty =}\ l_0 &= \tau_\infty\tau_1\\
\struck{\tau_\infty\tau_0 =}\ l_1 &= \tau_0\tau_\infty\\
l_\infty &= \tau_1\tau_0
\end{aligned}
\right.
\]
\note{les deux premières lignes commencent par une écriture biffée d'un
hachurage, lue $\tau_1\tau_\infty =$ et $\tau_\infty\tau_0 =$ ; en face, à
droite, « \textbf{NB} » souligné, sans suite}

$\mathcal{E}_\tau$ peut être considéré comme le groupe engendré par les $\tau_i$,
avec les seules relations (16) --- \ill{} \ill{} \ill{} $\tau_i \to -1$ de
$\mathcal{E}_\tau$ dans $\pm 1$, dont le noyau est $\pi$ engendré par les $l_i$
définis par (17)…
\marginal{\ill{} relations \ill{}}

Revenons à une description plus symétrique \uncertain{des extensions de}
$\mathcal{E}$. On désigne par $\hat{\mathcal{E}}$ le complété profini de $\mathcal{E}$,
extension de $\hat\Gamma_{0,3} = \Gamma_{0,3} = \mathfrak{S}_3 \times
\mathbb{Z}/2$ par $\hat\pi$, et par $\hat{\hat{\mathcal{E}}}$ le groupe des
\uncertain{automorphismes} \struck{\uncertain{extérieurs}} \uncertain{lacets}
de $\hat\pi$, extension de $\hat{\hat\Gamma} =
\operatorname{Autext}_{\mathrm{lac}}(\hat\pi)$ par $\hat\pi$ :
\[
(18)\quad 1 \longrightarrow \hat\pi \longrightarrow \hat{\hat{\mathcal{E}}}
\longrightarrow \hat{\hat\Gamma} \longrightarrow 1,
\qquad \hat{\hat\Gamma} \subset \operatorname{Autext}_{\mathrm{lac}}(\hat\pi)
\]
\note{sous $\hat{\hat\Gamma}$ de la suite (18), un petit signe vertical, lu
comme une inclusion (ou une égalité) avec
$\operatorname{Autext}_{\mathrm{lac}}(\hat\pi)$ écrit en dessous. L'indice
« lac » est lu sous « Autext » ici et sur toute la suite}

\page{23}

qui contient $\hat{\mathcal{E}}$ comme sous-extension. On désigne par
$\boldsymbol{\Gamma}$ le commutant de $\mathfrak{S}_3$ dans
$\hat{\hat\Gamma}$ (\ill{} \ill{}), et par $\mathbf{E}$ l'extension
correspondante de $\boldsymbol{\Gamma}$ par $\hat\pi$. Notons déjà que
\[
\left\{
\begin{aligned}
&\Gamma \cap \boldsymbol{\Gamma} = \{1, \tau\}\\
&\boldsymbol{\Gamma} \cap \mathfrak{S}_3 = \{1\}
\end{aligned}
\right.
\]
\note{$\boldsymbol{\Gamma}$ rend un $\Gamma$ à double jambage. Dans la
première ligne le second $\Gamma$ est encadré d'un trait appuyé qui recouvre
un $\mathfrak{S}_3$ ; la lecture $\Gamma \cap \boldsymbol{\Gamma}$ est
incertaine. À droite, une rature}

\marginal{$\mathfrak{S}_3 \times \hat{\hat\Gamma}^{!}$ \quad $\cup$ \quad
\textbf{NB} $\hat{\hat\Gamma} \supset \mathfrak{S}_3 \times
\boldsymbol{\Gamma}$ --- j'ignore si elles \uncertain{coïncident}, \ill{} une
égalité, i.e.\ si $\hat{\hat\Gamma}^{!} \to \boldsymbol{\Gamma}$ \ill{}
\uncertain{est} une égalité}
\note{la note marginale est écrite en oblique ; elle introduit un
$\hat{\hat\Gamma}^{!}$ que la page ne définit pas autrement et que la suite
utilise ($\hat{\hat{\mathcal{E}}}^{!}$)}

On veut donner d'abord une description plus symétrique de
$\hat{\hat{\mathcal{E}}}^{!}$, sans faire jouer un rôle particulier à
$L_0$ = \uncertain{sous-groupe engendré} par $l_0$
($L_0 = k_0(\hat{\mathbb{Z}})$). Un $u \in \hat{\hat{\mathcal{E}}}$ est défini par
un multiplicateur
\[
(19)\quad p = \chi(u) \in \hat{\mathbb{Z}}^{*}
\]
et des
\[
(20)\quad g_i \in \hat\pi \qquad i = 0, 1, \infty
\]
\struck{\ill{}}
\[
(21)\quad u(l_i) = \operatorname{int}(g_i^{-1})\, l_i^{\,p}
\]
\note{l'exposant $-1$ sur $g_i$, ici et dans toute la suite, est ajouté après
coup d'un crayon plus noir}
les données (19), (20) étant liées par la relation
\[
(22)\quad (\operatorname{int}(g_\infty^{-1})\, l_\infty^{\,p})\,
(\operatorname{int}(g_1^{-1})\, l_1^{\,p})\,
(\operatorname{int}(g_0^{-1})\, l_0^{\,p}) = 1,
\]
et par le fait que les $u(l_i)$ définis par (21) engendrent $\hat\pi$
topologiquement. Pour $u$ fixé, $p$

\page{24}
\note{p. 11 de l'auteur}

est unique, et $g_i$ défini \uncertain{à} multiplication \uncertain{à gauche}
près par un élt.\ quelconque de $L_i$, soit $l_i^{\alpha_i}$
($\alpha_i \in \hat{\mathbb{Z}}$). Les \ill{} de $u$ : $\operatorname{int}(g)
\circ u$ \uncertain{gardent} $p$, et \uncertain{remplacent} les $g_i$ par les
$gg_i$.

Si $v$ est donné par $q, (h_i)$, on a
\[
vu(l_i) = v(\operatorname{int}(g_i^{-1})\, l_i^{\,p})
= \operatorname{int}(v(g_i^{-1}))\, v(l_i)^{p}
\]
où
\[
v(l_i) = \operatorname{int}(h_i^{-1})\, l_i^{\,q} \quad \text{donc} \quad
v(l_i)^{p} = \operatorname{int}(h_i^{-1})\, l_i^{\,pq},
\]
d'où
\[
(23)\quad vu(l_i) = \operatorname{int}(v(g_i^{-1})\, h_i^{-1})\, l_i^{\,pq}
\]
\note{au-dessus de l'argument de $\operatorname{int}$, une première écriture
biffée, lue $(h_i\, v(g_i))^{-1}$, reliée par une accolade}

Si on désigne par $S\hat{\hat{\mathcal{E}}}^{!}$ l'extension de
$\hat{\hat{\mathcal{E}}}^{!}$ par $\hat{\mathbb{Z}}^{\{0,1,\infty\}}$ définie
\struck{\ill{}} \uncertain{par le} groupe des \struck{couples} triples d'un
$u \in \hat{\hat{\mathcal{E}}}^{!}$ et d'un \ill{} (21) \ill{} des $g_i$ qui
explicite \ill{} ($p$ étant le multiplicateur), on a sur
$S\hat{\hat{\mathcal{E}}}^{!}$ \uncertain{trois} fonctions à valeurs dans $\hat\pi$
\[
(u, (g_i)) \overset{\gamma_i}{\longmapsto} g_i ,
\]
et on aura
\[
(24)\quad \gamma_i(\tilde v\tilde u) =
\struck{\tilde v(\gamma_i(\tilde u))\cdot\gamma_i(\tilde v)}\
\tilde v(\gamma_i(\tilde u))\,\gamma_i(\tilde v)
\]
\note{deux premières écritures du membre de droite sont biffées, l'une par
hachurage, l'autre d'un trait ; la troisième, écrite à droite et reliée par
un trait courbe, est celle qui est gardée}
i.e.\ $(\gamma_i)$ est un 1-cocycle sur le groupe $S\hat{\hat{\mathcal{E}}}^{!}$
\struck{\ill{}} à valeurs dans le groupe \struck{\ill{} \ill{}}
$\hat\pi^{\{0,1,\infty\}}$.
Le sous-groupe \struck{\ill{}} de $S\hat{\hat{\mathcal{E}}}^{!}$ noyau de
$S\hat{\hat{\mathcal{E}}}^{!} \to \operatorname{Autext}^{!}_{\mathrm{lac}}(\hat\pi)$
s'identifie au groupe produit
\[
\Bigl(\hat\pi \times \prod_{i \in \{0,1,\infty\}} L_i\Bigr) \simeq
\hat\pi \times \hat{\mathbb{Z}}^{\{0,1,\infty\}}
\]

\marginal{\textbf{NB} $S\hat{\hat{\mathcal{E}}}^{!}$ a été construit comme ensemble
des systèmes $(p, g_0, g_1, g_\infty) \in \hat{\mathbb{Z}}^{*} \times
\hat\pi^{\{0,1,\infty\}}$ satisfaisant la relation (22)}

\page{25}

en associant à $(g, (\alpha_i)_{i \in \{0,1,\infty\}})$ le système
\[
\struck{\ill{}}\ \struck{\ill{}}\ (g\, l_i^{\alpha_i})
\]
\struck{dans} $p = 1$, $(g_\cdot) = \struck{\ill{}}\,(g\, l_i^{\alpha_i})$.
\note{« $p=1$, $(g_\cdot) =$ » ouvre la ligne ; les deux groupes biffés qui
suivent sont illisibles, et la parenthèse $(g\,l_i^{\alpha_i})$ est ce qui
reste}

Il faudrait expliciter les opérations des éléments $\sigma_i$ ou
$\varepsilon_i$ ($\in S\hat{\hat{\mathcal{E}}}^{!}$), ou $\tilde\varepsilon$
\ill{} $\in S\hat{\hat{\mathcal{E}}}^{!}$, sur ces groupes, et en particulier
\ill{} que \struck{\ill{}} les $\gamma_i$ sont « covariants » par $\rho$,
i.e.
\[
\gamma_i(\rho\tilde u\rho^{-1}) = \gamma_{i+1}(\tilde u)
\]
\struck{i.e.\ \ill{} $\tilde u$ est défini par}

\marginal{\textbf{NB} Il y a lieu d'introduire aussi le groupe
$S\hat{\mathcal{E}}$, extension de $\mathcal{E}$ par $\hat{\mathbb{Z}}^{\{0,1,\infty\}}$
(\uncertain{celle-ci} étant d'ailleurs une extension de $\Gamma$ par
\struck{\ill{}} \ill{} \ill{} $\hat{\mathbb{Z}}^{\{0,1,\infty\}}$) et
d'introduire \ill{} les relèvements des $\sigma_i$, $\rho$ dans
$S\hat{\mathcal{E}}$}
\page{26}
\note{p. 12 de l'auteur}

L'extension $\mathcal{E} = \mathcal{E}_{0,3}$ de $\Gamma = \Gamma_{0,3} \simeq
\mathfrak{S}_{\{0,1,\infty\}} \times \mathbb{Z}/2\mathbb{Z}$ par
\[
\pi = \pi_{0,3} = \{ l_0, l_1, l_\infty \mid l_\infty l_1 l_0 = 1 \}
= \pi_1(\mathbb{P}^1_{\mathbb{C}} \setminus \{0,1,\infty\}, \ill{}) .
\]
\note{sous $\mathfrak{S}_{\{0,1,\infty\}}$, une accolade avec « $=\mathfrak{S}_3$ » ;
sous $\mathbb{Z}/2\mathbb{Z}$, « engendré par $\tau : z \mapsto \bar z$ » ;
sous $\mathbb{P}^1_{\mathbb{C}} \setminus \{0,1,\infty\}$, une accolade avec
« $U = U_{0,3}$ » (lecture incertaine). Le point base, après la virgule, est
illisible}

Les classes de conjugaison \struck{\uncertain{au-dessus}} de scindages partiels
maximaux correspondent aux points fixes de $\dot\rho, \dot\sigma_0,
\dot\sigma_1, \dot\sigma_\infty$ opérant sur $X_{0,3}$, et
(\uncertain{éventuellement}) aux comp.\ connexes des ens.\ d'invariants de
$\dot\rho\tau, \dot\sigma_0\tau, \dot\sigma_1\tau, \dot\sigma_\infty\tau$. Les
pts fixes de $\dot\rho$ \uncertain{sont} (savoir $\varphi(\mathfrak{S}_3)$)
$P$ et $\overline{P}$ (dont le stabilisateur est d'ordre 6, formé de $1,
\dot\rho, \dot\rho^2, \dot\sigma_0\tau, \dot\sigma_1\tau,
\dot\sigma_\infty\tau$). On a étudié le relèvement correspondant (de
$\varphi(\mathfrak{S}_3)$) : $P$ ; on trouve \uncertain{entre eux} des
relèvements (\uncertain{$\pi$-conjugués}) relatifs à $\overline{P}$
s'obtiennent en \struck{conj} prenant un élément de $\mathcal{E}$
\struck{\ill{}} qui \uncertain{soit} au-dessus d'un élément de $\Gamma$
envoyant $P$ en $\overline{P}$ (savoir $\dot\sigma_0, \dot\sigma_1,
\dot\sigma_\infty, \tau, \tau\rho, \tau\rho^2$) et en conjuguant par cet
élément. On peut donc prendre \struck{\ill{}} un des éléments $\sigma_0,
\sigma_1, \sigma_\infty$ (ou $\varepsilon_0, \varepsilon_1,
\varepsilon_\infty$) ou $\tau_0, \tau_1, \tau_\infty$ ou \struck{\ill{}} les
produits des précédents par \struck{\ill{}} $\rho$ ou $\rho^2$….

Le \add{seul} \struck{les} pt\struck{s} fixe de $\sigma_0$ est $2$, dont le
stabilisateur est le groupe $\simeq \mathbb{Z}/2 \times \mathbb{Z}/2$
engendré par $\sigma_0$ et $\tau$. On a un relèvement de ce sous-groupe

\page{27}

dans $\mathcal{E}$, qui \uncertain{par} choix convenable \ill{} prenant le
quart de grand cercle qui joint $P$ au milieu $2$ du \struck{segment} l'arc
$(1, \infty)$ de grand cercle opposé à $0$, ce qui définit un
\uncertain{isomorphisme} \ill{} \ill{} universel de $U$ en $P$, avec celui
de $U$ en $\overline{P}$ sur lequel le groupe \uncertain{précédent} opère
de façon canonique, \ill{} \uncertain{aussi} \ill{} \ill{} structure sur
$\widetilde{U}(P)$, \ill{} il s'envoie dans $\mathcal{E}$.
\note{« $\widetilde{U}(P)$ » est écrit sur une rature appuyée ; l'argument
entre parenthèses est lu $P$ sans certitude}

On trouve que
\[
\dot\sigma_0 \longmapsto \sigma_0, \qquad \struck{t}\ \tau \longmapsto \tau_0
\]
est le \uncertain{scindage} correspondant. \uncertain{De} $\equiv$ pour le
pt fixe $-1$ de $\dot\sigma_1$, ou $\frac12$ de $\dot\sigma_\infty$, on
trouve des scindages partiels sur les ss-groupes de \struck{\ill{}}
stabilité $\Gamma_{-1} \subset \Gamma$ et $\Gamma_{\frac12} \subset
\Gamma$. Je \uncertain{note} \ill{} \ill{} \ill{}
\note{les indices des deux sous-groupes de stabilité sont lus $-1$ et
$\frac12$ d'après les points fixes nommés juste avant ; sur la page le
second ressemble à un $2$}
\[
\struck{F}\ 2 = 0' , \quad -1 = 1' ; \quad \tfrac12 = \infty' ,
\]
d'où des groupes d'isotropie $\Gamma_{0'} = \{\dot\sigma_0, \tau\}$,
$\Gamma_{1'} = \{\dot\sigma_1, \tau\}$ et $\Gamma_{\infty'} =
\{\dot\sigma_\infty, \tau\}$, qui se remontent \uncertain{chacun} en
$\sigma_0$ et $\tau_0$, \struck{\ill{}} en $\sigma_1$ et $\tau_1$, en
$\sigma_\infty$ et $\tau_\infty$.

Les composantes connexes de $U^{\tau}$ sont les trois arcs de grand cercle
entre $0, 1, \infty$, \struck{\uncertain{contiennent}} \uncertain{contenant}
chacun un \struck{\ill{}} et un seul des pts $0', 1', \infty'$, et les
\struck{classes} classes de conjugaison

\page{28}
\note{p. 13 de l'auteur}

correspondantes de relèvements (\add{relèvements} du $\{1, \dot\tau\}$) sont
resp.\ les relèvements en $\tau_0$, en $\tau_1$, et en $\tau_\infty$ --- qui
ne donnent rien de nouveau. Reste à considérer $\tau\dot\sigma_1$ (et ses
conjugués \uncertain{par} $\rho$, $\tau\dot\sigma_\infty$,
$\tau\dot\sigma_0$), on a $\dot\sigma_1(2) = 1/2$, donc
\[
(\tau\dot\sigma_1)(2) = \tfrac12
\]
et
\[
U^{\tau\dot\sigma_1} = \{ z \in U \mid z\bar z = 1 \} = \mathbb{U} - \{1\},
\]
où $\mathbb{U} = \{ z \in \mathbb{C} \mid |z| = 1 \}$. Cet ensemble contient
$P$ et $\overline{P}$, et ne donne donc pas de nouveaux relèvements p.r.\ à
celui déjà donné, avec $\tilde\sigma_1$ relevant $\dot\sigma_1\tau$.

Je \uncertain{voudrais} quand même donner aussi une description de $\mathcal{E}$
\uncertain{et} \uncertain{notamment} par gén.\ et relations, en partant d'un
syst.\ de gén.\ de $\Gamma = \mathfrak{S}_3 \times \mathbb{Z}/2$, qu'on
relève de façon judicieuse. On prendra les (\uncertain{gén.}) générateurs
\[
\struck{\ill{}} = \rho\,\struck{\ill{}},\ \tilde{\dot\sigma}_0,\
\tilde{\dot\sigma}_1,\ \tilde{\dot\sigma}_\infty
\qquad (\tilde{\dot\sigma}_0 = \dot\sigma_0\tau,\ \tilde{\dot\sigma}_1 =
\dot\sigma_1\tau,\ \tilde{\dot\sigma}_\infty = \dot\sigma_\infty\tau)
\]
(donc $\rho'^{6} = 1$, $\rho'^{3} = \tau$, $\rho'^{2} = \rho^{2} =
\rho^{-1}$) avec les relations
\[
\left\{
\begin{aligned}
&\tilde{\dot\sigma}_0^{\,2} = \tilde{\dot\sigma}_1^{\,2} =
\tilde{\dot\sigma}_{\uncertain{3}}^{\,2} = 1\\
&\tilde{\dot\sigma}_0\tilde{\dot\sigma}_1 =
\tilde{\dot\sigma}_1\tilde{\dot\sigma}_\infty =
\tilde{\dot\sigma}_\infty\tilde{\dot\sigma}_0 = \rho'^{\,\ill{}}\\
&\rho'^{6} = 1\\
&\rho'^{3} \text{ commute aux } \dot\sigma_i
\end{aligned}
\right.
\]
\note{le nouveau générateur est écrit sur une rature appuyée, au début de la
liste ; il est noté ici $\rho'$, car les exposants lus « 16 » et « 13 » sont
vraisemblablement $\rho'^{6}$ et $\rho'^{3}$, l'apostrophe tracée comme un
1. Tout le bloc final (la liste des générateurs et les relations) est biffé
de longs traits en diagonale. Dans la première relation, le troisième indice
est lu « 3 » là où l'on attend $\infty$ ; l'exposant de $\rho'$ dans la
deuxième est surchargé}

\page{29}

\struck{qui décrivent $\mathfrak{S}_3 \times \mathbb{Z}/2$, on en relève
ces générateurs en des élts de $\mathcal{E}$, de façon que $l_\infty, l_1, l_0$
s'expriment comme des mots en ces éléments, et il suffit alors d'ajouter
aux relations précédentes… la relation $l_\infty l_1 l_0$ ---
\uncertain{la} relation correspondante}
\note{ce paragraphe, qui achève la phrase biffée au bas de la page 28, est
lui-même biffé de quatre traits obliques}

\uncertain{Or} \ill{} \ill{} $\dot\sigma_0, \dot\sigma_1, \dot\sigma_\infty,
\dot\rho, \tau$, avec les relations
\[
\left\{
\begin{aligned}
&\dot\sigma_0^{2} = \dot\sigma_1^{2} = \dot\sigma_\infty^{2} = 1\\
&\dot\sigma_0\dot\sigma_1 = \dot\sigma_1\dot\sigma_\infty =
\dot\sigma_\infty\dot\sigma_0 = \dot\rho\\
&\dot\rho^{3} = 1\\
&\tau^{2} = 1\\
&\tau \text{ commute à } \dot\sigma_0, \dot\sigma_1, \dot\sigma_\infty
\quad (\text{donc aussi à } \dot\rho)
\end{aligned}
\right.
\]
On les relève en \struck{$\tau_0, \tau_1, \tau_\infty$} $\sigma_0, \sigma_1,
\sigma_\infty, \rho, \tau_0, \tau_1, \tau_\infty$ (ces trois derniers
au-dessus de $\tau$), et on trouve les relations
\[
\left\{
\begin{aligned}
&\sigma_0^{2} = \sigma_1^{2} = \sigma_\infty^{2} = \tau_0^{2}
= [\tau_1^{2} = \tau_\infty^{2}] = 1\\
&\rho^{3} = 1\\
&\rho\tau_0\rho^{-1} = \tau_1, \quad \rho\tau_1\rho^{-1} = \tau_\infty,
\quad \rho\tau_\infty\rho^{-1} = \tau_0\\
&\tau_0\sigma_0 = \sigma_0\tau_0, \quad \tau_1\sigma_1 = \sigma_1\tau_1,
\quad \tau_\infty\sigma_\infty = \sigma_\infty\tau_\infty\\
&\sigma_0\sigma_1 = \tau_\infty\tau_1\rho, \quad \sigma_1\sigma_\infty =
\tau_0\tau_\infty\rho, \quad \sigma_\infty\sigma_0 = \tau_1\tau_0\rho
\end{aligned}
\right.
\]

\marginal{\textbf{NB} Il faut ajouter les relations $\rho\sigma_0\rho^{-1} =
\sigma_1$, $\rho\sigma_1\rho^{-1} = \sigma_\infty$, $\rho\sigma_\infty\rho^{-1}
= \sigma_0$, conséquences des relations ci-contre (\uncertain{définissant}
les éléments $\sigma_0, \sigma_1, \sigma_\infty$)}
\note{la note marginale est écrite en oblique dans la marge gauche et reliée
par une flèche à la ligne $\tau_0\sigma_0 = \sigma_0\tau_0$ ; la lecture de la
parenthèse finale est incertaine}

\marginal{\textbf{NB} \uncertain{On trouve} $l_0 = \tau_\infty\tau_1$,
$\struck{l_1 = \tau_0\tau_\infty}$, $l_1 = \tau_0\tau_\infty$, $l_\infty =
\tau_1\tau_0$ et $l_\infty l_1 l_0 = 1$ \ill{} \ill{} \uncertain{Cela}
indique que $\tau_0, \tau_1, \tau_\infty$ engendrent un sous-groupe engendré
par $l_0, l_1, l_\infty$}
\note{note entre crochets à droite du système, reliée par une flèche montante
à la relation $\rho\tau_\infty\rho^{-1} = \tau_0$ ; la dernière phrase est
lue avec peine}

\struck{\uncertain{Question} : Est-ce \uncertain{que} $\rho \mapsto$
\uncertain{composantes} $\rho$ en fonction des $\sigma_0, \sigma_1,
\sigma_\infty, \tau_0, \tau_1, \tau_\infty$\,? --- l'une des relations, et
\ill{} \ill{} dernier groupe de relations, et \uncertain{substituer} dans
\uncertain{les} \uncertain{autres}.}
\note{ce paragraphe, qui ferme la page, est biffé de deux traits obliques ;
un segment « et \ill{} plus » y est en outre barré à part}

\page{30}
\note{p. 14 de l'auteur}

On peut prendre aussi comme \uncertain{seuls} gén.\ $\rho, \sigma_0, \tau_0$
(dont se déduisent $\sigma_1, \tau_1$, comme $\rho\sigma_0\rho^{-1},
\rho\tau_0\rho^{-1}$, $\sigma_\infty, \tau_\infty$ comme
$\rho^{2}\sigma_0\rho^{-2}, \rho^{2}\tau_0\rho^{-2}$) et les relations
\[
\left\{
\begin{aligned}
&\text{a)}\ \sigma_0^{2} = \tau_0^{2} = 1, \quad \tau_0\sigma_0 =
\sigma_0\tau_0 \quad (\text{i.e.\ } (\underbrace{\sigma_0\tau_0}_{\tilde\sigma_0})^{2} = 1)\\
&\text{b)}\ \rho^{3} = 1\\
&\underbrace{\sigma_0\rho\sigma_0^{-1}\rho^{-1}}_{\uncertain{\sigma_1}}
= \underbrace{\rho^{2}\tau_0\rho^{-2}}_{\tau_\infty}\,
\underbrace{\rho\tau_0\rho^{-1}}_{\tau_1}\,\rho
\quad \text{i.e.} \quad \sigma_0\rho\sigma_0^{-1}\rho^{-1} =
\rho^{-1}\tau_0\rho\tau_0\\
&\rho\sigma_0\rho^{-1}\sigma_0 = \tau_0\rho^{-1}\tau_0\rho
\quad \text{\uncertain{soit}} \quad \struck{\operatorname{int}(\rho)\sigma_0}\
[\rho, \sigma_0] = [\tau_0, \rho^{-1}]\\
&\text{\uncertain{soit}} \quad [\rho, \sigma_0][\rho^{-1}, \tau_0] = 1
\end{aligned}
\right.
\]
\note{au-dessus de « $[\rho,\sigma_0] = [\tau_0,\rho^{-1}]$ », des mots
biffés (lus « i.e. \ill{} \ill{} ») ; au-dessus de « $\tau_0\rho^{-1}\tau_0
\rho$ », « \uncertain{ou encore} ». L'accolade sous $\sigma_0\rho\sigma_0^{-1}
\rho^{-1}$ porte un $\sigma_1$, que la seule écriture $\rho\sigma_0\rho^{-1}$
aurait donné}

(C'est donc un quotient du produit libre
\[
(\mathbb{Z}/2\mathbb{Z} \times \mathbb{Z}/2\mathbb{Z}) * \mathbb{Z}/3\mathbb{Z},
\]
par le sous-groupe invariant décrit par la relation $[\rho, \sigma_0][\rho^{-1},
\tau_0] = 1$.

\textbf{NB} On trouve une symétrie dans les groupes, qui échange $\sigma_0$ et
$\tau_0$, et échange $\rho$ et $\rho^{-1}$. [donc qui échange \ill{}
$\sigma_1$ et $\tau_\infty$, $\sigma_\infty$ et $\tau_1$] --- mais cette
symétrie n'est pas compatible avec la structure d'extension, comme on voit
de suite (\uncertain{en passant au quotient} $\mathfrak{S}_3 \times
\mathbb{Z}/2$, elle devrait \uncertain{conserver} les \uncertain{cartes}\,!).

\page{31}

Il faudrait encore expliciter l'extension $S\Gamma = S\Gamma_{0,3}$ de
$\Gamma$ par $\mathbb{Z}^{\{0,1,\infty\}}$, ce qui nous donnera (par produits
fibrés) une description de
\[
S\mathcal{E} = S\Gamma \times_{\Gamma} \mathcal{E} .
\]

\page{32}
\note{p. 15 de l'auteur}

\subsection{(2) Le groupe $S\hat{\hat{\mathcal{E}}}$}
\note{le titre est de sa main, souligné ; le « 2 » est cerclé}

Par le calcul, il s'explicite comme l'ens.\ des systèmes
\[
(1)\quad (p, \sigma, g_0, g_1, g_\infty) = \tilde u, \qquad
p \in \hat{\mathbb{Z}}^{*},\ \sigma \in \mathfrak{S}_3,\ g_0, g_1, g_\infty
\in \hat\pi,
\]
satisfaisant la condition
\[
(2)\quad \operatorname{int}(g_\infty)(l_{\sigma(\infty)}^{\,p}) \cdot
\operatorname{int}(g_1)(l_{\sigma(1)}^{\,p}) \cdot
\operatorname{int}(g_0)(l_{\sigma(0)}^{\,p}) = 1
\]
\note{dans le premier facteur, $g_\infty$ est écrit sur un $g_0$ surchargé.
L'exposant $-1$ des pages 23--24 a disparu des $g_i$ sous
$\operatorname{int}$}
et tels que les trois facteurs de ce produit soient un syst.\ de générateurs
de $\hat\pi$ [avec la relation génératrice (2)]. Le morphisme $u$ défini
par ce \struck{morphi} système est
\[
(3)\quad
\left\{
\begin{aligned}
u(l_0) &= \operatorname{int}(g_0)\, l_{\sigma(0)}^{\,p}\\
u(l_1) &= \operatorname{int}(g_1)\, l_{\sigma(1)}^{\,p}\\
u(l_\infty) &= \operatorname{int}(g_\infty)\, l_{\sigma(\infty)}^{\,p}
\end{aligned}
\right.
\qquad \text{ce qui explicite l'homomorphisme}
\]
\note{dans la première ligne de (3), le signe $=$ est écrit sur une flèche
$\mapsto$}
\[
(4)\quad S\hat{\hat{\mathcal{E}}} \xrightarrow{\ p\ } \hat{\hat{\mathcal{E}}} =
\operatorname{Aut}_{\mathrm{lac}}(\hat\pi) .
\]
L'homomorphisme canonique
\[
(5)\quad i : \hat\pi \times \hat{\mathbb{Z}}^{\{0,1,\infty\}} \longrightarrow
S\hat{\hat{\mathcal{E}}}
\]
s'explicite par
\[
(6)\quad i(g; \alpha_0, \alpha_1, \alpha_\infty) = (1, 1, g\,l_0^{\alpha_0},
g\,l_1^{\alpha_1}, g\,l_\infty^{\alpha_\infty})
\]
donc
\[
(7)\quad p\,i(g; \alpha_0, \alpha_1, \alpha_\infty) = \operatorname{int} g .
\]

\marginal{On retrouve que (4) est surjectif, et a comme noyau
$i(\hat\pi)$}
\note{note écrite en oblique dans la marge gauche, face aux formules (5)--(7) ;
elle dit « $i(\hat\pi)$ » là où l'on attendrait l'image de $i$ tout entière}

La multiplication dans $S\hat{\hat{\mathcal{E}}}$ est donnée par
\[
(8)\quad \underbrace{(p, \sigma, g_0, g_1, g_\infty)}_{\tilde u}
\underbrace{(p', \sigma', g'_0, g'_1, g'_\infty)}_{\tilde u'} =
(pp', \sigma\sigma', u(g'_0)\,g_{\sigma'(0)}, u(g'_1)\,g_{\sigma'(1)},
u(g'_\infty)\,g_{\sigma'(\infty)})
\]
\note{les indices des $g$ du membre de droite sont lus $\sigma'(0)$,
$\sigma'(1)$, $\sigma'(\infty)$ ; le dernier est au bord de la feuille}

Cette formule n'est pas, du point de vue du calcul,

\page{33}

entièrement explicite et satisfaisante, puisque le deuxième membre fait
intervenir les quantités $u(g'_0), u(g'_1), u(g'_\infty)$, qui ne sont donc
pas exprimées directement en termes des \add{seules} données
$(p, \sigma, g_0, g_1, g_\infty)$, $(p', \sigma', g'_0, g'_1, g'_\infty)$ ;
\struck{\ill{}} pour \uncertain{que} $g'_0, g'_1, g'_\infty$ soient
eux-mêmes donnés explicitement comme des mots (ou limites de tels) en les
$l_0, l_1, l_\infty$, mais ils s'explicitent ainsi.
\note{la dernière proposition, « mais ils s'explicitent ainsi », est
ajoutée au-dessus de la ligne ; son point d'insertion n'est pas sûr}

Cependant quand $\tilde u$ est de la forme $i(g, \alpha_0, \alpha_1,
\alpha_\infty)$, on trouve (en tenant compte des \uncertain{facteurs})
\[
(9)\quad i(g, \alpha_0, \alpha_1, \alpha_\infty)\,(p, \sigma, g_0, g_1,
g_\infty) = (p, \sigma, g\,g_0\,l_{\sigma(0)}^{\alpha_{\sigma(0)}},
g\,g_1\,l_{\sigma(1)}^{\alpha_{\sigma(1)}},
g\,g_\infty\,l_{\sigma(\infty)}^{\alpha_{\sigma(\infty)}})
\]
\note{devant le premier $g\,g_0$, une surcharge biffée ; les indices sont
lus avec peine}
qui se simplifie pour $\sigma = 1$ en
\[
(9')\quad i(g, \alpha_0, \alpha_1, \alpha_\infty)\,(p, 1, g_0, g_1,
g_\infty) = (p, 1, g\,g_0\,l_0^{\alpha_0}, g\,g_1\,l_1^{\alpha_1},
g\,g_\infty\,l_\infty^{\alpha_\infty}) .
\]
Par ailleurs, l'opération de $\hat{\hat{\mathcal{E}}}$ sur $i(\pi \times
\mathbb{Z}^{\{0,1,\infty\}})$ par automorphismes intérieurs s'explicite
aisément, et ne dépend que de $u$ \struck{\ill{}} :
\[
(10)\quad \tilde u\, i(g; \alpha_0, \alpha_1, \alpha_\infty)\,\tilde u^{-1}
= i(u(g); p\,\alpha_{\sigma^{-1}(0)}, p\,\alpha_{\sigma^{-1}(1)},
p\,\alpha_{\sigma^{-1}(\infty)})
\]
où $\tilde u = (p, \sigma; g_0, g_1, g_\infty)$ --- mais sur cette formule, on
peut faire les mêmes commentaires que plus haut, sur la nécessité
d'expliciter des \uncertain{termes} $u(g)$ de cette formule, pour $g$ non
exprimé comme un mot p.r.\ aux $l_0, l_1, l_\infty$ (ou limite de tels mots).
\note{dans (10), l'exposant $-1$ de $\tilde u^{-1}$ est ajouté d'une encre
plus noire sur un signe surchargé}

\page{34}
\note{p. 16 de l'auteur}

On va regarder de plus près la partie $S\hat{\mathcal{E}}$ de
$S\hat{\hat{\mathcal{E}}}$, image inverse de $\Gamma = \mathfrak{S}_3 \times
\mathbb{Z}/2 = \hat\Gamma \subset \hat{\hat\Gamma} =
\operatorname{Autext}_{\mathrm{lac}}\hat\pi$.
On va regarder les éléments \struck{\ill{}} $\rho$, $\sigma_0, \sigma_1,
\sigma_\infty$, $\tau_0, \tau_1, \tau_\infty$ de $\hat{\hat{\mathcal{E}}}$,
\struck{\uncertain{notamment}} (\textbf{NB} ils sont \uncertain{même} dans
$\mathcal{E}$) et voir comment ils se relèvent en $S\hat{\mathcal{E}}$ (et même, en
$S\mathcal{E}$). En fait, par $\dot\rho \mapsto \rho$, $\dot\sigma_0 \mapsto
\tilde\sigma_0 = \sigma_0\tau_0$, on a un relèvement multiplicatif de
$\Gamma^{*}_{P} = \{(\sigma, \operatorname{sg}(\sigma)) \mid \sigma \in
\mathfrak{S}_3\}$ dans $\mathcal{E}$, qui aimerait remonter à $S\mathcal{E}$ --- et de
même pour $\dot\sigma_0 \mapsto \sigma_0$, $\tau \mapsto \tau_0$ du
ss-groupe de stabilité de $0$ dans $\Gamma$ (et, par conjugaison, des pour
$1, \infty$). On trouve en fait un homomorphisme canonique
\note{« $\Gamma^{*}_{P}$ » : l'astérisque est tracé comme une petite étoile,
l'indice est un $P$ ; le mot « multiplicatif » est ajouté au-dessus de la
ligne. Le « $= \hat\Gamma$ » sous $\mathfrak{S}_3 \times \mathbb{Z}/2$ est
écrit en dessous, relié par un signe d'égalité vertical}
\[
(11)\quad \varphi : \Gamma^{*}_{P} \longrightarrow S\mathcal{E}
\]
\struck{qui remonte} section partielle de $S\mathcal{E}$ sur $\Gamma$, qui remonte
$\Gamma^{*} \to \mathcal{E}$ déjà considéré, par
\[
(12)\quad
\begin{aligned}
\dot\rho = (\dot\rho, 1) &\longmapsto (1, \rho, 1, 1, 1)
\overset{\text{déf}}{=} \tilde\rho\\
\dot\sigma'_i = (\dot\sigma_i, -1) &\longmapsto (-1, \dot\sigma_i, 1, 1, 1)
\overset{\text{déf}}{=} \tilde\sigma'_i \qquad \bigl(i \in \{0, 1,
\infty\}\bigr)
\end{aligned}
\]
\note{dans la première ligne de (12), le second terme du quintuplet est un
$\rho$ sans point ; dans la seconde, un $\dot\sigma_i$ pointé}

\struck{On} peut dire aussi que le sous-\struck{groupe} \add{ens.}\ de
$S\hat{\hat{\mathcal{E}}}$ formé des $(p, \sigma, g_0, g_1, g_\infty)$ pour lesquels
$g_0, g_1, g_\infty$ sont égaux à $1$ est un \emph{sous-groupe}
\note{« ens. » est écrit au-dessus de « groupe » biffé}

\page{35}

(évident sur (8)), qui \struck{\uncertain{est}} \uncertain{peut se} décrire
comme l'ens.\ des $(p, \sigma)$ ($p \in \hat{\mathbb{Z}}^{*}$, $\sigma \in
\mathfrak{S}_3$) tels que l'on ait
\[
(13)\quad l_{\sigma(0)}^{\,p}\, l_{\sigma(1)}^{\,p}\, l_{\sigma(\infty)}^{\,p} = 1 .
\]
Quand on se borne à $S\mathcal{E}$ \ill{} lieu de $S\hat{\hat{\mathcal{E}}}$, de sorte
qu'on doit avoir $p \in \mathbb{Z}^{*} = \{\pm 1\}$, \struck{on trouve que
les \ill{}} i.e.\ l'ens.\ des $(p, \sigma)$ \struck{\uncertain{tels}} \add{à}
la condition (13) \uncertain{devient} \struck{\ill{}}\,$\times \mathfrak{S}_3
\simeq \mathfrak{S}_3 \times \mathbb{Z}_2 \simeq \Gamma$), \uncertain{il}
n'est autre que $\Gamma_P$ : on voit que $(p, \sigma)$ satisfait (13) ssi
$(\sigma, p) \in \Gamma_P$. Ainsi, $\tilde\Gamma_P \subset S\mathcal{E} \subset
S\hat{\hat{\mathcal{E}}}$ est caractérisé comme le sous-groupe de $S\mathcal{E}$ des
auto\add{morphismes} de $\pi$ \struck{\ill{}} qui invarient $L_0, L_1,
L_\infty$, \uncertain{affectés} des « rectificateurs »
\struck{multiplicateurs} $g_0 = 1$, $g_1 = 1$, $g_\infty = 1$.
\note{le « $\Gamma_P$ » est écrit sans l'astérisque de la page 34 ; le groupe
biffé devant « $\times\,\mathfrak{S}_3$ » est illisible sous le hachurage.
Dans « \ill{} lieu », le mot manquant est sans doute « au », mais il est
tracé comme « a »}

Sans doute il est vrai que le \uncertain{même} résultat est valable dans
$S\hat{\hat{\mathcal{E}}}$, i.e.\ que la relation (13) implique forcément $p \in
\{\pm 1\}$. Je \uncertain{pense} \uncertain{qu'on} \uncertain{peut}
[\uncertain{quitte} : corriger par un $\tilde\sigma_i$, ou $\sigma \in
\Gamma_P$\,)] revient au $\equiv$ de dire que la relation
\[
(14)\quad l_\infty^{\,p}\, l_1^{\,p}\, l_0^{\,p} = 1
\]
(pour un $p \in \hat{\mathbb{Z}}^{*}$ fixé) implique $p = 1$.

Exercice\,!
\note{la parenthèse entre crochets est écrite sous la ligne et reliée à «
Je » ; la phrase qui commence par « Je » est lue avec peine}

\page{36}
\note{p. 17 de l'auteur}

Comment remonter les $2$-\uncertain{groupes} de Sylow $\simeq \mathbb{Z}/2
\times \mathbb{Z}/2$ de $\Gamma$, engendrés resp.\ par $\{\sigma_0, \tau\}$,
$\{\sigma_1, \tau\}$, $\{\sigma_\infty, \tau\}$\,? Eh bien justement, ils
ne se remontent \emph{pas} --- on sait bien que les extensions locales de
$\Gamma^{+}_0$ (p.\,ex.) par $\mathbb{Z}_0$ (quotient de
$\mathbb{Z}^{\{0,1,\infty\}}$ par $\mathbb{Z}^{\{1,\infty\}}$) n'est
\emph{pas} scindée, i.e.\ le générateur $\dot\sigma_0$ de $\Gamma^{+}_0$ ne
se remonte \emph{pas} en un élément d'ordre $2$ de $\mathcal{E}$ qui centralise
(ou ce qui revient au $\equiv$, normalise) $L_0$.
\note{le numérateur de « $2$-\uncertain{groupes} de Sylow » est surchargé ;
« de Sylow » est ajouté au-dessus de la ligne. « $\Gamma^{+}_0$ » : la
première fois, l'exposant est une croix appuyée, lue $+$ ; « (p.\,ex.) » est
ajouté au-dessus}

On pourrait s'amuser à écrire $S\Gamma$ et $S\mathcal{E}$ en générateurs et
relations --- je ne vais pas m'y attarder. Par contre, je vais donner des
formules pour l'action de $\operatorname{int}(\tilde\rho)$ et des
$\operatorname{int}(\tilde\sigma_i)$. On trouve
\[
\begin{aligned}
(15)&\quad (\operatorname{int}\tilde\rho)(p, \sigma, g_0, g_1, g_\infty) =
(p, \dot\rho\sigma\dot\rho^{-1}, \rho(g_\infty), \rho(g_0), \rho(g_1))\\
(16)&\quad \operatorname{int}(\tilde\sigma_0)(p, \sigma, g_0, g_1, g_\infty) =
(p, \dot\sigma_0\sigma\dot\sigma_0^{-1}, \sigma_0(g_0), \sigma_0(g_\infty),
\sigma_0(g_1))
\end{aligned}
\]
\note{dans (15), le second terme est surchargé ; il est lu
$\dot\rho\sigma\dot\rho^{-1}$ par analogie avec (16)}

On voit donc : a) Pour que $\tilde u = (p, \sigma, g_0, g_1, g_\infty)$
commute à $\rho$, il f.\ et s.\ que \struck{$g_0, g_1, g_\infty$ soient}
$\sigma$ commute à $\rho$, i.e.\ $\sigma \in \{1, \rho, \rho^{2}\}$,
\emph{et} qu'on ait
\[
(17)\quad g_1 = \rho(g_0), \quad g_\infty = \rho(g_1), \quad g_0 =
\rho(g_\infty) .
\]
Ce qui signifie aussi que $g_1, g_\infty$ se déduisent de $g_0$ par
\[
g_1 = \rho(g_0), \quad g_\infty = \rho(g_1) = \rho^{2}(g_0) =
\rho^{-1}(g_0)
\]

\page{37}

(qui impliquent déjà $g_0 = \rho(g_\infty)$ dans (17)). Les $\tilde u$ en
question correspondent donc aux systèmes $(p, \sigma, g_0) \in
\hat{\mathbb{Z}}^{*} \times \mathfrak{S}_3^{+} \times \hat\pi$ tels que l'on
ait la relation qui correspond à (2) :
\[
(18)\quad \operatorname{int}(\rho^{2}(g_0))(l_{\sigma(\infty)}^{\,p}) \cdot
\operatorname{int}(\rho(g_0))(l_{\sigma(1)}^{\,p}) \cdot
\operatorname{int}(g_0)(l_{\sigma(0)}^{\,p}) = 1
\]
(et que les trois facteurs engendrent $\hat\pi$…)
\note{dans (18), l'exposant du deuxième facteur porte un astérisque ajouté
au-dessus du $p$ ; « $\mathfrak{S}_3^{+}$ » est lu tel quel}

b) Pour que $\tilde u$ commute à $\tilde\sigma_0$, il f.\ et s.\ que
$\sigma$ commute à $\dot\sigma_0$ (i.e.\ $\sigma \in \{1, \dot\sigma_0\}$),
et qu'on ait
\[
g_0 = \sigma_0(g_0), \quad g_\infty = \sigma_0(g_1) \quad (\text{donc } g_1 =
\sigma_0(g_\infty))
\]
\struck{\ill{}} \struck{\ill{}}

\marginal{\uncertain{variante} intéressante pour la suite}
\note{note écrite en oblique dans la marge gauche, face au début de b)}

À vrai dire, nous nous intéressons plutôt à la commutation dans
$\hat{\hat\Gamma}$, i.e.\ la commutation \uncertain{modulo} automorphismes
intérieurs.
\note{deux lignes biffées au-dessus de « À vrai dire », illisibles sous les
ratures ; la seconde commence par « \ill{} $\tilde u$ commute »}
\end{document}
